% Source pin (SHA-256): pifo-elegant-v4.md=356ea38146ced437e20200db04e2ce220c73f9ac9fa892b746f783b8a2d49ab8.
\subsection{Proper flush decompositions}
\label{sec:decompositions}

For induction on the alphabet, we expose the first queue that makes a
nontrivial routing decision.

\begin{lemma}[Proper top decomposition]
\label{lem:decomposition}
If \(|A|>1\), every valid stateless scheduler is online-equivalent to a
scheduler of the form
\begin{equation}
 M=(P,r,(M_B)_{B\in P}),
 \label{eq:decomposition-shape}
\end{equation}
where \(P\) is a proper partition of \(A\) into active root children, \(r\) is
the total preorder induced by root ranks, and \(M_B\) is the child scheduler
restricted to types in \(B\).  Its flush transformer satisfies
\begin{align}
 \restrict{F_M(w)}{B}
   &=F_{M_B}(\restrict{w}{B})
    =F_M(\restrict{w}{B}),
    \label{eq:block-projection}\\
 \proj{P}(F_M(w))
   &=\proj{P}(\sort_r(w)).
    \label{eq:root-pattern}
\end{align}
\end{lemma}

\begin{proof}
Starting at the root, follow the unique active child used by all types for as
long as one exists.  Each skipped queue contains only that child value and
therefore forwards exactly one service per successful pop; removing such a
node preserves every online trace.  The walk ends either at the first node with
at least two active children or at a leaf.  In the former case, partition types
by the chosen child.  Every block is nonempty and, because at least two are
active, proper.

In the leaf case, replace the leaf by an internal root having the old leaf rank
for each type and one FIFO singleton child per type.  The new root queue mirrors
the old leaf queue occurrence-for-occurrence, while a singleton child has no
choice to make.  Since \(|A|>1\), the singleton partition is proper.  This
establishes \autoref{eq:decomposition-shape} without changing online behavior.

During a flush, child \(B\) receives all pushes in
\(\restrict{w}{B}\) before it receives service, and the root calls it exactly
\(|\restrict{w}{B}|\) times.  Its complete output is therefore
\(F_{M_B}(\restrict{w}{B})\), proving the first equality of
\autoref{eq:block-projection}.  On a word confined to \(B\), all root references
select \(B\), which proves the second.  Finally the root drains its reference
occurrences by the stable \((r,\arr)\) order.  Replacing their source types by
their child blocks yields \autoref{eq:root-pattern}.
\end{proof}

We call \((P,r)\) a \emph{decomposition} of a word transformer \(F\) when
\autoref{eq:block-projection} and \autoref{eq:root-pattern} hold.  Two consequences will be used
without further comment:
\begin{align}
  F(v)&=F_{M_B}(v) &&\text{if every letter of \(v\) lies in \(B\)},
  \label{eq:block-autonomy-simple}\\
  \restrict{F(w)}{B}&=F(\restrict{w}{B}) &&\text{for arbitrary \(w\)}.
  \label{eq:block-filter-simple}
\end{align}
The first is \emph{block autonomy}; the second says that deleting other blocks
commutes with the transformer.

\subsection{Combining two decompositions}
\label{sec:meet}

Suppose \((P,r)\) and \((Q,s)\) decompose one common flush transformer \(F\).
Their meet partition is the set of nonempty incidence cells
\begin{equation}
 R=P\wedge Q
  =\{B\cap C\ne\varnothing:B\in P,\ C\in Q\}.
 \label{eq:meet-partition}
\end{equation}
We will construct a total preorder that agrees with \(r\) across \(P\)-blocks
and with \(s\) across \(Q\)-blocks.  We do \emph{not} need to claim that the
resulting pair \((R,t)\) itself decomposes \(F\).

First, every cell inherits autonomy.  For \(D=B\cap C\), repeated use of
\autoref{eq:block-filter-simple} gives
\begin{equation}
\begin{aligned}
 \restrict{F(w)}{D}
 &=\restrict{(\restrict{F(w)}{B})}{C}
  =\restrict{F(\restrict{w}{B})}{C}\\
 &=F(\restrict{(\restrict{w}{B})}{C})
  =F(\restrict{w}{D}).
\end{aligned}
\label{eq:cell-autonomy}
\end{equation}

For types in distinct \(R\)-cells, prescribe the effective comparison
\(e(x,y)\) recovered by \lemref{lem:two-type}.  If \(P\) separates the pair, its
two-type restriction is chosen at the \(P\)-root, and therefore
\begin{equation}
 e(x,y)=\cmp_r(x,y).
 \label{eq:e-from-r}
\end{equation}
Likewise, if \(Q\) separates it,
\begin{equation}
 e(x,y)=\cmp_s(x,y).
 \label{eq:e-from-s}
\end{equation}
The prescriptions agree when both partitions separate the pair because both
are recovered from the same two flushes of \(F\).

\begin{figure}[t]
\centering
\begin{tikzpicture}[
  cell/.style={minimum width=2.0cm,minimum height=1.2cm,draw},
  pt/.style={circle,fill=pifogray,text=white,inner sep=2.1pt,font=\small},
  every node/.style={font=\small}]
  \draw[step=2cm] (0,0) grid (4,2.4);
  \node[pt] (x) at (1,1.8) {$x$};
  \node[pt] (y) at (3,1.8) {$y$};
  \node[pt] (z) at (1,.6) {$z$};
  \node[above] at (1,2.4) {$C_1$};
  \node[above] at (3,2.4) {$C_2$};
  \node[left] at (0,1.8) {$B_1$};
  \node[left] at (0,.6) {$B_2$};
  \node[draw=pifored,very thick,rounded corners,fit=(x)(y),inner sep=10pt] {};
  \node[draw=pifoblue,very thick,rounded corners,fit=(x)(z),inner sep=10pt] {};
  \node[pifored,anchor=west] at (4.18,1.78) {$P$-block};
  \node[pifoblue,anchor=east] at (-.18,.32) {$Q$-block};
  \node[align=left,anchor=west] at (5.35,.62)
    {rows expose $\{x,y\}\mid\{z\}$;\\
     columns expose $\{x,z\}\mid\{y\}$};
\end{tikzpicture}
\caption{The only nontrivial three-cell incidence shape.  Three meet cells
form an L when neither all \(P\)-coordinates nor all \(Q\)-coordinates are
distinct.  Its two visible decompositions form the diamond used below.}
\label{fig:meet-diamond}
\end{figure}

\begin{lemma}[Three-cell diamond]
\label{lem:diamond}
On one representative type from each of any three distinct meet cells, the
pairwise comparisons \(e\) form a total preorder.
\end{lemma}

\begin{proof}
View the cells as occupied positions in the \(P\times Q\) incidence grid.  If
the three \(P\)-coordinates are distinct, all comparisons come from the genuine
preorder \(r\) by \autoref{eq:e-from-r}.  If their \(Q\)-coordinates are distinct,
they all come from \(s\) by \autoref{eq:e-from-s}.  Otherwise the three positions
form an L as in \autoref{fig:meet-diamond}; naming its corner \(x\), the restricted
transformer has both decompositions
\begin{equation}
  \{x,y\}\mid\{z\},
  \qquad
  \{x,z\}\mid\{y\}.
  \label{eq:l-decompositions}
\end{equation}

Assume that the non-strict relation induced by \(e\) is not transitive.  Since
each pair is totally compared, there are distinct \(a,b,c\) with
\begin{equation}
  a\le_e b,
  \qquad b\le_e c,
  \qquad c<_e a.
  \label{eq:bad-triple}
\end{equation}
Push one occurrence of each in the cycle order \(a,b,c\), then flush.  The
arrival stamps are distinct, so every queue computation has a unique least
\((\rank,\arr)\) key.  The cycle inequalities determine every strict or weak
root comparison, and the fixed arrival order resolves any remaining rank tie;
no separate tie case split is needed.

Consider in full the L orientation whose grouped pairs are \(\{b,c\}\) and
\(\{c,a\}\).  In the decomposition
\(\{b,c\}\mid\{a\}\), separated-pair comparisons give root ranks
\(\rho_a\le\rho_b\) and \(\rho_c<\rho_a\).  Thus
\(\rho_c<\rho_b\), and the root drains references in the order
\(c,a,b\): if \(\rho_a=\rho_b\), arrival puts \(a\) first.  The resulting
child slots are
\[
  \{b,c\},\ a,\ \{b,c\}.
\]
Within the pair child, \(b\) arrived before \(c\) and \(b\le_e c\), so the
binary recovery table says that its output is \(b,c\).  The complete output is
therefore \(b,a,c\).

In the decomposition \(\{c,a\}\mid\{b\}\), separated comparisons give
\(\sigma_a\le\sigma_b\le\sigma_c\).  With arrival order \(a,b,c\), the root
drains references in that same order, producing slots
\[
  \{c,a\},\ b,\ \{c,a\}.
\]
The pair child outputs \(c,a\), since \(c<_e a\), and the complete output is
\(c,b,a\).  The first letters differ, contradicting that both decompositions
compute the same \(F(abc)\).

Cyclic renaming covers the other two placements of the two grouped edges in
the comparison cycle.  Exchanging the roles of the two decompositions covers
the three reflected L orientations.  These exhaust the incidence shape, so no
failure of transitivity is possible.  Reflexivity and totality are already
built into the pair comparisons; hence the restriction is a total preorder.
\end{proof}

\begin{figure}[t]
\centering
\begin{tikzpicture}[
  panel/.style={draw,rounded corners,minimum width=6.1cm,minimum height=2.85cm,
                align=left,inner sep=8pt},
  title/.style={font=\bfseries}]
  \node[panel] (l) at (-3.35,0) {};
  \node[title,anchor=north] at (l.north) {$\{b,c\}\mid\{a\}$};
  \node[align=left,anchor=north west] at ($(l.north west)+(.18,-.48)$)
    {$\rho_c<\rho_a\le\rho_b$\\
     root sources: $c,a,b$\\
     slots: $[\{b,c\}], [a], [\{b,c\}]$\\
     inner output: $b,c$\\
     final output: $\mathbf{b,a,c}$};
  \node[panel] (r) at (3.35,0) {};
  \node[title,anchor=north] at (r.north) {$\{c,a\}\mid\{b\}$};
  \node[align=left,anchor=north west] at ($(r.north west)+(.18,-.48)$)
    {$\sigma_a\le\sigma_b\le\sigma_c$\\
     root sources: $a,b,c$\\
     slots: $[\{c,a\}], [b], [\{c,a\}]$\\
     inner output: $c,a$\\
     final output: $\mathbf{c,b,a}$};
  \draw[-{Stealth},very thick,pifored] ($(l.south)+(0,-.2)$) --
    node[below,font=\small] {contradict a common flush transformer}
    ($(r.south)+(0,-.2)$);
\end{tikzpicture}
\caption{The representative diamond computation for
\(a\le_e b\le_e c<_e a\), using one common arrival order \(a,b,c\).
Internal references choose child slots; the selected pair child may emit a
different source occurrence.}
\label{fig:diamond-computation}
\end{figure}
